\section{算出来：手术、$S$ 矩阵与 Verlinde 公式}

\subsection{要算什么}

到现在为止，我们有了一个理论：三维 Chern--Simons 理论，量子态是全纯截面。
但“理论”必须能算出数。这一讲算两类东西：
\begin{itemize}[leftmargin=*]
  \item 闭三维流形的配分函数 $Z(M)$；
  \item 带穿刺点曲面的 Hilbert 空间维数，即融合系数 $N_{ij}{}^k$。
\end{itemize}
算完之后，把纽结放进去，就得到 Jones 多项式。

\subsection{黎曼球上的三点函数}

先看最简单的曲面：$S^2$ 上三个标记点，表示分别为 $R_i,R_j,R_k$。对应的 Hilbert 空间维数记为
\begin{equation}
N_{ijk}=\dim\cH_{S^2;\{i,j,k\}}. \tag{7.1}
\end{equation}
这个数就是融合系数。$N_{ijk}\ne0$ 当且仅当 $R_k$ 出现在 $R_i\otimes R_j$ 的分解里。
在最简单的情形（大 $k$ 极限），$N_{ijk}$ 就是经典张量积的重数；有限 $k$ 时，它是它的一个子集
（只有可积表示才计入）。

\subsection{环面、$S$ 与 $T$}

亏格 $1$ 的情形给出另一个基本数据。环面 $\Sigma=T^2$ 上的 Hilbert 空间维数是 $k+1$（$SU(2)$ 时），
有一组自然基 $|j\rangle$，$j=0,\frac12,\dots,\frac k2$。模群 $SL(2,\ZZ)$ 作用在这个空间上，
由两个生成元给出：
\begin{equation}
T=\mathrm{diag}\left(e^{2\pi\ii(h_j-c/24)}\right),\qquad
S_{jj'}=\sqrt{\frac{2}{k+2}}\,\sin\frac{\pi(2j+1)(2j'+1)}{k+2}. \tag{7.2}
\end{equation}
$S$ 与 $T$ 满足 $S^2=(ST)^3=1$（在 Witten 的 $T$ 约定下；若把 $T$ 换成拓扑自旋
$\mathrm{diag}(e^{2\pi\ii h_j})$，则 $(ST)^3=e^{2\pi\ii c/8}$，差的正是标架反常）。

$S$ 的物理意义是“把环面沿一个圈切开再重新粘合”的变换。它也可以由共形块直接算出：
$S_{jj'}$ 是环面上的两个基态之间的内积，在模变换下交换了两个圈。

\subsection{手术}

任意一个闭三维流形都可以由 $S^3$ 沿一个带标架的链环做手术得到：挖掉每个链环分量的管状邻域
（实心环），再按标架换一种方式粘回去。每一步手术在量子理论里对应一个操作：
在切开处插入一个完备态，粘合时做求和。

具体地，沿一个纽结做手术时，插入的态是 $\sum_\lambda S_{0\lambda}\,|\lambda\rangle$ 型的组合，
其中 $S_{0\lambda}$ 是真空与 $\lambda$ 之间的 $S$ 矩阵元。于是
\begin{equation}
Z(M)=\sum_{\lambda_1,\dots,\lambda_r}
\left(\prod_{i}S_{0\lambda_i}\right)\;Z\!\left(S^3;\,\text{带标架的链环标号 }\lambda_i\right). \tag{7.3}
\end{equation}
这个公式把任意流形的配分函数化归为 $S^3$ 上带标号链环的配分函数，而后者可以用辫群与
共形块算出来。

\subsection{Verlinde 公式的推导}

现在考虑一个具体的比较：把 $S^2\times S^1$ 沿一个圈做手术，可以得到 $S^3$。
一方面，$S^2\times S^1$ 的 Hilbert 空间已知（就是 $S^2$ 的态空间）；另一方面，
手术之后应该得到 $S^3$ 的配分函数。两边相等，给出一个关于 $S$ 与 $N$ 的恒等式。

Witten 的推导（原文 (4.51)）给出
\begin{equation}
\sum_k S_i{}^k\,N_{kl}{}^m=\frac{S_{il}\,S_i{}^m}{S_{0i}}. \tag{7.4}
\end{equation}
这正是“$S$ 对角化融合规则”的数学表述。令 $v_i$ 为融合代数的基，$v_iv_j=\sum_k N_{ij}{}^kv_k$；
则在新基
\[
w_i=S_{0i}\sum_m S_i{}^m\,v_m
\]
下，乘法变成
\begin{equation}
w_iw_j=\delta_{ij}\,w_j, \tag{7.5}
\end{equation}
即 $w_i$ 是幂等元。于是所有融合系数都可以由 $S$ 算出：
\[
N_{ij}{}^k=\sum_l\frac{S_{il}S_{jl}S_{kl}}{S_{0l}}.
\]
这就是 Verlinde 公式。

\subsection{这个推导用了什么}

值得注意的是，上面这条推导链条只用到了三样东西：
\begin{enumerate}[leftmargin=*]
  \item 手术公式 (7.3)；
  \item $N_{ijk}$ 的定义 (7.1)；
  \item 态空间基的正交归一性。
\end{enumerate}
它\textbf{没有用到共形块的动力学}（KZ 方程、Ward 恒等式）。换句话说，Verlinde 公式的
“物理推导”实际上是拓扑场论自洽性的推论；把它变成纯数学定理，需要证明共形块空间确实满足
这些自洽性（TUY 1989、Faltings 1994 做的是这件事）。

\subsection{标架反常}

配分函数还依赖标架。改变标架时，
\begin{equation}
Z\longmapsto Z\cdot\exp\left(2\pi\ii\,\frac{n-m}{24}\,c\right), \tag{7.6}
\end{equation}
其中 $n-m$ 是标架改变对应的整数，$c$ 是 (6.3) 的中心荷。这个相位不影响 $|Z|$，
但影响带相位的量（例如 Jones 多项式在 $q$ 的单位根处的取值）。第八讲会再讨论它。

\subsection{环面配分函数与模不变性}

配分函数不只在闭三维流形上有定义；把时空取成 $\Sigma\times S^1$，得到的是一个二维配分函数
（对 $\Sigma$ 上的场求和），它作为 $S^1$ 的周长（模参数 $\tau$）的函数，满足模不变性。
对 $\Sigma=T^2$ 且无 Wilson 线的情形，
\[
Z_{T^2\times S^1}(\tau)=\sum_{j} |\chi_j(\tau)|^2 ,
\]
其中 $\chi_j$ 是可积表示的特征标，求和遍历可积表示。这个量在 $\tau\to-1/\tau$ 下不变，
因为 $S$ 矩阵是酉的、且特征标按 $S$ 变换。

这个公式是“模不变性”的最简形式，也是 Verlinde 公式的来源之一：
如果在环面上插入 Wilson 线（即标记点），配分函数变成 $\sum_j \chi_j\bar\chi_{j'}$ 的组合，
而融合系数正是这些组合的系数。换句话说，Verlinde 公式可以从“环面配分函数的模不变性”
推出来，这正是 Verlinde 原始论证的思路。

本文的推导走的是另一条路（手术 + 三维态空间），两条路给出同一个公式。
这种“多路验证”是有理共形场论的典型特征：一个公式常常有物理、代数、几何三种证明。

\subsection{一个可核对的数}

取 $G=SU(2)$、$k=2$。此时 $S_{00}=\sqrt{2/4}\sin(\pi/4)=1/2$，而 $S^3$ 的配分函数正是 $S_{00}$，
所以
\[
Z(S^3)=\tfrac12\qquad(k=2).
\]
一般地 $Z(S^3)=\sqrt{2/(k+2)}\sin(\pi/(k+2))$；$k=1$ 时它是 $1/\sqrt2$。
这个数可以用 (7.4) 与 $N$ 的定义独立核对，读者不妨自己算一遍。它是检验整套机器是否自洽的
一个最小例子。

\subsection{$S$ 矩阵从哪来}

(7.2) 的公式看起来像从天而降。其实它有一个干净的来源：环面态空间上的模变换。

取环面 $\Sigma=T^2$，记它的两个圈为 $a,b$。环面上的共形块（也就是三维理论在该曲面上的态）
由不可积表示 $|j\rangle$ 标号。模群 $SL(2,\ZZ)$ 由
\[
S=\begin{pmatrix}0&-1\\1&0\end{pmatrix},\qquad
T=\begin{pmatrix}1&1\\0&1\end{pmatrix}
\]
生成，它作用在态空间上。物理上，$T$ 是“沿 $a$ 圈做 Dehn 扭转”，
而 $S$ 是“交换 $a$ 与 $b$”。用共形场论的语言，$T$ 作用是对角矩阵
$\mathrm{diag}(e^{2\pi\ii(h_j-c/24)})$，因为扭转一圈给出共形权与中心荷的相位；
$S$ 作用则给出 (7.2) 的矩阵元。

(7.2) 的显式形式可以用 Weyl 特征标公式推出：$S_{jj'}$ 是特征标在模变换下的核，
对 $SU(2)$ 它就是 $\sqrt{2/(k+2)}\sin(\pi(2j+1)(2j'+1)/(k+2))$。
读者可以验证 $S^2=1$ 与 $(ST)^3=1$（在 Witten 的 $T$ 约定下）；这两个等式不是巧合，
它们就是模群关系的量子化版本。若把 $T$ 换成拓扑自旋，则 $(ST)^3$ 多出 $e^{2\pi\ii c/8}$，
这正是 (7.6) 的标架反常。

\subsection{辫群与 Jones 多项式}

最后一步是把链环写成辫子的闭合。任何定向链环都可以表示为一个辫子 $\beta\in B_n$ 的闭合
$\hat\beta$。在共形场论里，辫群作用在 $n$ 个标记点的共形块空间上：$B_n$ 的生成元
$\sigma_i$ 对应“把第 $i$ 与第 $i+1$ 个标记点交换”，而 KZ 方程给出的单值性正是这个作用。
闭合辫子对应于取迹：
\[
Z(S^3;\hat\beta)=\Tr_{\cH}\left(\rho(\beta)\right).
\]
对 $SU(2)$、定义表示，这个表示可以约化到 Temperley--Lieb 代数，其迹给出 Jones 多项式。
这就是 Jones 多项式“来自三维”的确切含义：它是共形块空间上辫群表示的一个迹。

\subsection{手术公式的推导}

把 (7.3) 写得更具体一些。设 $M$ 由 $S^3$ 沿链环 $L=L_1\cup\cdots\cup L_r$ 做手术得到，
标架为 $f_i$。手术的每一步可以分解为：
\begin{enumerate}[leftmargin=*]
  \item 挖去 $L_i$ 的管状邻域 $N_i\cong D^2\times S^1$，剩下一个边界为 $\coprod_i T^2_i$ 的流形 $M_0$；
  \item 在 $M_0$ 的每个边界环面上，态空间有一组基 $|\lambda_i\rangle$（由可积表示标号）；
  \item 实心环 $D^2\times S^1$ 的态是 $\sum_\lambda S_{0\lambda}|\lambda\rangle$
        （这是“真空在手术下的展开”，$S_{0\lambda}$ 是真空与 $\lambda$ 的 $S$ 矩阵元）；
  \item 粘合时做内积，得到求和。
\end{enumerate}
于是
\[
Z(M)=\sum_{\lambda_1,\dots,\lambda_r}
\left(\prod_i S_{0\lambda_i}\right)\,
Z(S^3;L\text{ 标号 }\lambda_1,\dots,\lambda_r).
\]
这个公式的严格版本需要证明：手术不依赖具体的链环表示（Kirby  calculus）。
Kirby 定理给出链环图之间的三种基本移动，而 $S$ 矩阵与融合系数满足的恒等式
（$S^2=(ST)^3=1$、Verlinde 公式）恰好保证 $Z(M)$ 在这些移动下不变。
这就是“为什么手术公式是自洽的”的答案：自洽性等价于模群关系。

\subsection{Heegaard 分裂与 WRT 不变量}

手术公式有一个等价的表述，叫 Heegaard 分裂。任何闭三维流形都可以写成两个柄体
沿一个曲面 $\Sigma$ 粘合：$M=M_1\cup_\Sigma M_2$，其中 $M_1,M_2$ 各是一个柄体。
在 Chern--Simons 理论里，每个柄体给 $\Sigma$ 上一个态（向量），粘合就是内积：
\[
Z(M)=\langle \psi_1,\psi_2\rangle .
\]
这个表述比手术公式更几何：它把三维流形的不变量化归为曲面上的态空间与两个向量。
而曲面上的态空间，正是第六讲的共形块空间。

Reshetikhin--Turaev 的构造可以看作这个表述的组合化：用一个模张量范畴（或量子群在单位根处的
表示范畴）替换“共形块”，用范畴的缝合运算替换“内积”，得到对链环与三维流形的组合不变量。
这个构造的优点是完全组合、可计算；缺点是它把“为什么”丢掉了——范畴从哪里来、
为什么与 Chern--Simons 一致，需要 Witten 的路径积分来解释。

对读者来说，两条路都值得知道：手术公式告诉你“怎么算”，Heegaard 分裂告诉你“为什么这样算”。
本文的贡献是用后者解释前者：三维理论在曲面上的态空间就是共形块，所以手术公式不是技巧，
而是理论的结构。

\subsection{一般群与 $S$ 的公式}

对一般的紧致单群 $G$ 与 level $k$，$S$ 矩阵由 Weyl 特征标公式给出：
\begin{equation}
S_{\lambda\mu}=\frac{1}{|W|}\sum_{w\in W}\varepsilon(w)\,
\exp\left(-\frac{2\pi\ii}{k+h^\vee}\langle w(\lambda+\rho),\mu+\rho\rangle\right), \tag{7.7}
\end{equation}
其中 $W$ 是 Weyl 群，$\varepsilon(w)$ 是符号，求和遍历 Weyl 群。
对 $SU(2)$，$W=\{1,-1\}$，把 $\lambda=2j$ 代入就回到 (7.2) 的正弦形式。
对 $SU(N)$，$S$ 矩阵可以用 Schur 多项式与单位根处的取值写出；
这些公式在 Verlinde 的原始论文与后来的综述里都有。

(7.7) 的一个重要推论是：$S$ 矩阵只依赖 level $k$ 与表示，不依赖曲面的亏格。
亏格只影响态空间的维数（由 Verlinde 公式给出），不影响融合规则。
这就是“有理共形场论”的结构：数据是有限维的、亏格无关的。

\subsection{三叶结：把数算出来}

最后用一个具体例子把整套机器串起来。三叶结是最简单的非平凡纽结，
它的 Jones 多项式是
\[
J(\text{三叶结})=-t^{-4}+t^{-3}+t^{-1}.
\]
在 Chern--Simons 的框架里，三叶结是 $2$ 股辫子 $\sigma_1^3$ 的闭合。
$SU(2)$、定义表示的辫群表示可以约化到 Temperley--Lieb 代数 $TL_n$，
其生成元 $e_i$ 满足
\[
e_i^2=\delta\,e_i,\qquad e_ie_{i\pm1}e_i=e_i,\qquad e_ie_j=e_je_i\ (|i-j|\ge2),
\]
其中 $\delta=-q-q^{-1}$，$q=e^{2\pi\ii/(k+2)}$。
辫子生成元 $\sigma_i$ 可以写成 $q^{1/2}(1-e_i)$ 的倍数；闭合取迹（Markov 迹）后，
得到的量就是 Jones 多项式在 $q$ 处的值。

取 $k=1$（此时 $q=e^{2\pi\ii/3}$），$TL_2$ 只有两个基，$e_1$ 是一个 $2\times2$ 矩阵；
$\sigma_1^3$ 的迹可以直接算出，得到 $J$ 在 $q$ 处的值。
把 $q$ 看作变量并解析延拓，就得到上面的多项式。

这个计算的意义在于：它不是“验证”，而是“定义”。在 Chern--Simons 框架里，
Jones 多项式的变量 $t$ 就是 $q^2$，它的单位根取值对应有限的 $k$；
多项式本身是这些取值的插值。这就是第一讲所说的“三维定义”的具体含义。

\subsection{手算 $k=2$ 的 $S$ 矩阵}

$k=2$ 时 $SU(2)$ 的可积表示是 $j=0,\frac12,1$，共三个。用 (7.2)：
\[
S_{jj'}=\sqrt{\frac{2}{4}}\,\sin\frac{\pi(2j+1)(2j'+1)}{4}
=\frac{1}{\sqrt2}\sin\frac{\pi(2j+1)(2j'+1)}{4}.
\]
逐项代入：
\[
S=\frac{1}{\sqrt2}
\begin{pmatrix}
\sin\frac{\pi}{4} & \sin\frac{2\pi}{4} & \sin\frac{3\pi}{4}\\[2pt]
\sin\frac{2\pi}{4} & \sin\frac{4\pi}{4} & \sin\frac{6\pi}{4}\\[2pt]
\sin\frac{3\pi}{4} & \sin\frac{6\pi}{4} & \sin\frac{9\pi}{4}
\end{pmatrix}
=
\frac12
\begin{pmatrix}
1 & \sqrt2 & 1\\
\sqrt2 & 0 & -\sqrt2\\
1 & -\sqrt2 & 1
\end{pmatrix}.
\]
于是 $S_{00}=\frac12$，与前面 $Z(S^3)=S_{00}=\frac12$ 一致；$S^2=1$ 可以逐项验证
（注意 $\sqrt2\cdot\sqrt2=2$），$(ST)^3=1$ 需要先写出 $T$，留给读者。

这个矩阵还有一个用途：用它算融合系数。例如
\[
N_{1/2,1/2,1}=\sum_{l=0,\frac12,1}\frac{S_{1/2,l}^2\,S_{1,l}}{S_{0,l}}
=\frac{\tfrac12\cdot\tfrac12}{\tfrac12}+\frac{\tfrac12\cdot\tfrac12}{\tfrac12}=1,
\]
第一项来自 $l=0$（$S_{1/2,0}^2=\tfrac12$，$S_{1,0}=\tfrac12$，$S_{0,0}=\tfrac12$），
第二项来自 $l=1$（$S_{1/2,1}^2=\tfrac12$，$S_{1,1}=\tfrac12$），$l=\frac12$ 项为零。
同理 $N_{1/2,1/2,0}=1$。于是 $1/2\otimes1/2=0\oplus1$，
与第六讲的融合规则一致。

\subsection*{练习}
\begin{enumerate}[leftmargin=*]
  \item 用 (7.2) 验证 $S^2=1$（$SU(2)$、$k=1,2$）。
  \item 计算 $Z(S^3)$ 在 $k=1,2,3$ 时的值，并验证它们等于 $S_{00}$。
  \item 用 (7.4) 与 $k=1$ 的 $S$ 矩阵验证 $N_{1/2,1/2,0}=1$、$N_{1/2,1/2,1/2}=0$。
\end{enumerate}

\begin{quote}
\small
\textbf{本讲要点。} 手术公式把 $Z(M)$ 化归为 $S^3$ 上的链环；$S$ 与 $T$ 由环面态空间给出 (7.2)；
(7.4) 是 Verlinde 公式的等价形式；标架反常 (7.6) 带 $c/24$；Jones 多项式是辫群表示在共形块上的迹。
\end{quote}
